\section{Introduction}

Discrete diffusion Large Language Models (dLLMs) have emerged as a compelling alternative to autoregressive generation, offering the potential for non-monotonic reasoning and parallel decoding. However, the standard absorbing-state framework—which enforces a rigid, monotonic transition from $\text{[MASK]}$ to fixed tokens—faces inherent limitations in fidelity. As highlighted by \citet{kang_parallelbench_2025}, the independent nature of parallel decoding often amplifies token-level inconsistencies. While recent studies have attempted to mitigate this via confidence-based remasking \citep{wang_remasking_2025} or by employing external guide models \citep{lee_effective_2025}. To bridge the gap between efficient parallel generation and high-fidelity reasoning, we align with the direction of generalizing discrete diffusion beyond absorbing states \citep{rutte_generalized_2025} and propose a comprehensive framework for Editable State Evolution.

Unlike prior work such as \citet{song2025seed}, we first design a novel \textbf{Error-Correcting Editable} decoding strategy, which introduces a dynamic paradigm controlled by dual probability thresholds. This paradigm encompasses two types of operations: direct decoding from mask to token, and editing from one token to another. This strategy enables the model to directly refine its own outputs during the generation process, thereby effectively addressing the local inconsistencies commonly encountered in parallel decoding. To cultivate this editing capability, our CPT and SFT phases expose the model to both masked positions and stochastic noise, incentivizing it to not only generate new content but also identify and rectify existing errors.

\textbf{Crucially, this architecture transforms the rigid trade-off between latency and fidelity into a flexible, user-configurable continuum.} By allowing the model to retroactively correct errors, we can aggressively lower the confidence threshold for the initial Mask-to-Token (M2T) phase without collapsing the generation quality. This insight gives rise to two distinct operating personas: a \textbf{\textit{Speedy Mode (S Mode)}}, which prioritizes high-throughput generation by accepting lower-confidence tokens and relying on subsequent Token-to-Token (T2T) passes for rectification; and a \textbf{\textit{Quality Mode (Q Mode)}}, which adheres to conservative thresholds to maximize reasoning rigor. This duality demonstrates that editability is not merely a mechanism for error repair, but a fundamental lever for accelerating parallel decoding.


To further elevate the model's capabilities, we integrate a \textbf{Reinforcement Learning (RL)} stage. While recent works such as SPG \citep{wang2025spg}, TraceRL \citep{wang2025revolutionizing} and ESPO \citep{ou_principled_2025} have demonstrated the potential of RL in improving dLLMs, applying policy gradients to block-autoregressive models remains challenging due to the intractability of sequence log-likelihoods. We circumvent this by adopting an ELBO-based Block-level Policy Optimization (EBPO) framework tailored for our editable setting. 

Notice that LLaDA2.1 extends its previous version (LLaDA2.0) by prioritizing decoding versatility over mere parameter scaling or benchmark peaking. By keeping the model size constant and minimal change of training data, we prove that our novel editing scheme enables lightning-fast execution with minimal overhead. This work serves as a proof-of-concept for a new dLLM paradigm that balances high-quality generation with extreme operational efficiency.

